\documentclass[11pt,a4paper]{article}
\usepackage[utf8]{inputenc}
\usepackage[spanish]{babel}
\usepackage{amsmath, amssymb}
\usepackage{geometry}
\geometry{top=2.5cm, bottom=2.5cm, left=2.5cm, right=2.5cm}

\begin{document}

\section*{Piense, dialogue y comparta}

\begin{itemize}
    \item[\textbf{1.}] Usted está trabajando en una planta de fabricación. Una de las máquinas de la planta utiliza un resorte. Para un nuevo proceso, sería deseable que se utilizaran resortes con diferentes constantes elásticas durante diferentes partes del proceso. Su jefe le pide consejo sobre cómo cambiar un resorte por otro rápidamente para que todo el proceso pueda tener lugar en una cantidad razonable de tiempo. Usted sugiere que, en lugar de cambiar resortes continuamente, usted cambiaría la constante elástica de uno de los resortes largos sujetándolo en diferentes posiciones para definir un nuevo extremo fijo del resorte. Entonces el resorte efectivo consiste solamente en las bobinas que están más allá de la abrazadera. Usted diseña un sistema que consta de un resorte largo con $N$ bobinas y una constante elástica $k$. Usted diseña un sistema de abrazadera que aislará parte del resorte, dejando $N'$ bobinas más allá de la abrazadera fija. (a) Escriba una expresión para la constante elástica $k'$ del extremo libre del resorte en términos de $k, N$ y $N'$. (b) El extremo libre del resorte es sujetado y jalado hacia afuera una distancia $x$. En el proceso, la mano que sostiene el extremo del resorte hace un trabajo $W$ en el resorte. Ahora, se regresa el resorte a su estado relajado y luego se fija en su punto central. El extremo libre del resorte relajado con abrazadera, se sujeta y se jala hacia fuera la misma distancia $x$. ¿Cuánto trabajo hace la mano en el resorte en este caso? ¿$W$? ¿$2W$? ¿$4W$? ¿Otro valor?
    \vspace*{9cm}

    \item[\textbf{2.}] \textbf{ACTIVIDAD} En la tabla de datos, vemos las distancias de frenado mínimas $d$ en función de la rapidez inicial $v$ de un auto. Trabaje en su grupo para responder lo siguiente. (a) Si duplica la rapidez inicial, ¿se necesita el doble de distancia para detener el auto? (b) Suponga que la distancia de frenado es proporcional a la rapidez del auto elevada a cierta potencia: $d \propto v^n$. Usar técnicas gráficas para determinar $n$. (c) ¿Por qué la distancia de frenado depende del valor particular de $n$ que encontró en (b)?
    \vspace*{8cm}
\end{itemize}

\newpage

\section*{Problemas}

\subsection*{SECCIÓN 8.1 Modelo de análisis: Sistema no aislado (energía)}

\begin{itemize}
    \item[\textbf{1.}] Una bola de masa $m$ cae desde una altura $h$ al piso. (a) Escriba la versión adecuada de la ecuación 8.2 para el sistema de la bola y la Tierra y utilícela para calcular la rapidez de la bola justo antes de que pegue en la Tierra. (b) Escriba la versión apropiada de la ecuación 8.2 para el sistema de la bola y utilícela para calcular la rapidez de la pelota justo antes de que pegue en la Tierra.
    \vspace*{6cm}
\end{itemize}

\subsection*{SECCIÓN 8.2 Modelo de análisis: Sistema aislado (energía)}

\begin{itemize}
    \item[\textbf{2.}] Una bola de cañón de $20.0 \text{ kg}$ se dispara desde un cañón con rapidez de boquilla de $1\,000 \text{ m/s}$ con un ángulo de $37.0^\circ$ con la horizontal. Una segunda bala de cañón se dispara con un ángulo de $90.0^\circ$. Aplique el modelo de sistema aislado para encontrar (a) la altura máxima que alcanza cada bola y (b) la energía mecánica total del sistema bola-Tierra en la altura máxima para cada bola. Sea $y = 0$ en el cañón.
    \vspace*{6cm}

    \item[\textbf{3.}] Un bloque de masa $m = 5.00 \text{ kg}$ se libera desde el punto \textcircled{A} y se desliza sobre la pista sin fricción que se muestra en la figura P8.3. Determine (a) la rapidez del bloque en los puntos \textcircled{B} y \textcircled{C} y (b) el trabajo neto hecho por la fuerza gravitacional en el bloque cuando se mueve de \textcircled{A} a \textcircled{C}.
    \vspace*{6cm}

    \item[\textbf{4.}] A las 11:00 a.m. del 7 de septiembre de 2001, más de 1 millón de escolares británicos saltaron hacia arriba y hacia abajo durante un minuto para simular un terremoto. (a) Encuentre la energía almacenada en los cuerpos de los niños que fue convertida en energía interna en el suelo y en sus cuerpos y se propagó en el suelo mediante ondas sísmicas durante el experimento. Suponga que los $1\,050\,000$ niños de masa promedio $36.0 \text{ kg}$ saltaron 12 veces cada uno, elevando sus centros de masa por $25.0 \text{ cm}$ cada vez y descansando brevemente entre un salto y el siguiente. (b) De la energía que se propaga en el suelo, la mayoría produce vibraciones de alta frecuencia ``microtemblores'' que fueron rápidamente amortiguadas y no viajaron mucho. Suponga $0.01\%$ de la energía total se transportó lejos por ondas sísmicas de largo alcance. La magnitud de un sismo en la escala de Richter está dada por: $M = \frac{\log E - 4.8}{1.5}$, donde $E$ es la energía de la onda sísmica en joules. De acuerdo con este modelo, ¿cuál fue la magnitud de sismo de demostración?
    \vspace*{6cm}
\end{itemize}

\newpage

\begin{itemize}
    \item[\textbf{5.}] Una barra ligera rígida mide $77.0 \text{ cm}$ de largo. Su extremo superior está articulado en un eje horizontal sin fricción. La barra cuelga recta hacia abajo en reposo con una pequeña bola de gran masa unida a su extremo inferior. Usted golpea la bola y súbitamente le da una velocidad horizontal de modo que logra un círculo completo. ¿Qué rapidez mínima se requiere en la parte más baja para hacer que la bola recorra lo alto del círculo?
    \vspace*{6cm}

    \item[\textbf{6.}] \textbf{Problema de repaso.} El sistema que se muestra en la figura P8.6 consiste en una cuerda ligera inextensible; poleas ligeras sin fricción; y bloques de igual masa. Observe que el bloque B está unido a una de las poleas. Inicialmente el sistema se mantiene en reposo de modo que los bloques están a la misma altura sobre el suelo. Después los bloques se liberan. Encuentre la rapidez del bloque A en el momento en que la separación vertical de los bloques es $h$.
    \vspace*{6cm}
\end{itemize}

\subsection*{SECCIÓN 8.3 Situaciones que incluyen fricción cinética}

\begin{itemize}
    \item[\textbf{7.}] Una caja de $10.0 \text{ kg}$ de masa se jala hacia arriba de un plano inclinado rugoso con una rapidez inicial de $1.50 \text{ m/s}$. La fuerza del jalón es $100 \text{ N}$ paralela al plano, que forma un ángulo de $20.0^\circ$ con la horizontal. El coeficiente de fricción es $0.400$ y la caja se jala $5.00 \text{ m}$. (a) ¿Cuánto trabajo hace la fuerza gravitacional en la caja? (b) Determine el aumento en energía interna del sistema caja-plano inclinado debido a fricción. (c) ¿Cuánto trabajo hace la fuerza de $100 \text{ N}$ en la caja? (d) ¿Cuál es el cambio en energía cinética de la caja? (e) ¿Cuál es la rapidez de la caja después de que se ha jalado $5.00 \text{ m}$?
    \vspace*{8cm}

    \item[\textbf{8.}] Una caja de $40.0 \text{ kg}$, inicialmente en reposo, se empuja $5.00 \text{ m}$ a lo largo de un suelo horizontal rugoso, con una fuerza constante horizontal aplicada de $130 \text{ N}$. El coeficiente de fricción entre la caja y el suelo es $0.300$. Encuentre: (a) el trabajo hecho por la fuerza aplicada, (b) el aumento en energía interna en el sistema caja-suelo como resultado de la fricción, (c) el trabajo hecho por la fuerza normal, (d) el trabajo hecho por la fuerza gravitacional, (e) el cambio en energía cinética de la caja y (f) la rapidez final de la caja.
    \vspace*{8cm}
\end{itemize}

\newpage

\begin{itemize}
    \item[\textbf{9.}] Se coloca un aro circular plano con radio de $0.500 \text{ m}$ sobre el piso. Una partícula $0.400 \text{ kg}$ se desliza alrededor del borde interior del aro. A la partícula se le da una rapidez inicial de $8.00 \text{ m/s}$. Después de una revolución, su rapidez cae a $6.00 \text{ m/s}$ debido a la fricción con el suelo rugoso de la pista. (a) Encuentre la energía transformada de mecánica a interna en el sistema partícula-aro-piso como resultado de la fricción en una revolución. (b) ¿Cuál es el número total de revoluciones que da la partícula antes de detenerse? Suponga que la fuerza de fricción permanece constante durante todo el movimiento.
    \vspace*{6cm}
\end{itemize}

\subsection*{SECCIÓN 8.4 Cambios en la energía mecánica para fuerzas no conservativas}

\begin{itemize}
    \item[\textbf{10.}] Como se muestra en la figura P8.10, una cuenta verde de masa $25 \text{ g}$ se desliza a lo largo de un alambre recto. La longitud del alambre desde el punto \textcircled{A} al punto \textcircled{B} es $0.600 \text{ m}$ y el punto \textcircled{A} está $0.200 \text{ m}$ más arriba que el punto \textcircled{B}. Una fuerza de rozamiento constante de magnitud $0.025 \text{ N}$ actúa en la cuenta. (a) Si la cuenta se libera a partir del reposo en el punto \textcircled{A}. (b) Una cuenta roja de masa $25 \text{ g}$ se desliza a lo largo de un alambre curvado, sujeto a una fuerza de fricción con la misma magnitud constante que la cuenta verde. Si las cuentas verde y roja se liberan al mismo tiempo a partir del reposo en el punto \textcircled{A}, ¿qué cuenta alcanza el punto \textcircled{B} con una rapidez mayor? Explique.
    \vspace*{6cm}

    \item[\textbf{11.}] En un tiempo $t_i$, la energía cinética de una partícula es $30.0 \text{ J}$ y la energía potencial del sistema al que pertenece es $10.0 \text{ J}$. En algún tiempo posterior $t_f$, la energía cinética de la partícula es $18.0 \text{ J}$. (a) Si solo fuerzas conservativas actúan en la partícula, ¿cuáles son la energía potencial y la energía total en el tiempo $t_f$? (b) Si la energía potencial en el tiempo $t_f$ es $5.00 \text{ J}$, ¿existen fuerzas no conservativas que actúan en la partícula? (c) Explique su respuesta del inciso (b).
    \vspace*{6cm}

    \item[\textbf{12.}] Un objeto de $1.50 \text{ kg}$ se mantiene $1.20 \text{ m}$ sobre un resorte vertical relajado sin masa con una constante elástica de $320 \text{ N/m}$. Se deja caer el objeto sobre el resorte. (a) ¿Cuánto comprime al resorte? (b) \textbf{¿Qué pasaría si?} Repita el inciso (a), pero esta vez suponga que una fuerza de resistencia al aire constante de $0.700 \text{ N}$ actúa en el objeto durante su movimiento. (c) \textbf{¿Qué pasaría si?} ¿Cuánto comprime el objeto al resorte si el mismo experimento se realiza en la Luna, donde $g = 1.63 \text{ m/s}^2$ y se desprecia la resistencia del aire?
    \vspace*{7cm}
\end{itemize}

\newpage

\begin{itemize}
    \item[\textbf{13.}] Una niña de masa $m$ parte del reposo y se desliza sin fricción desde una altura $h$ a lo largo de un tobogán al lado de una piscina (figura P8.13). Se lanzó desde una altura $h/5$ en el aire sobre la piscina. Queremos encontrar la altura máxima que alcanza arriba del agua en su movimiento como proyectil. (a) ¿Es el sistema niña-Tierra aislado o no aislado? ¿Por qué? (b) ¿Hay una fuerza no conservativa que actúa dentro del sistema? (c) Defina la configuración del sistema cuando la niña está en el nivel del agua que tiene cero energía potencial gravitacional. Exprese la energía total del sistema cuando la niña está en la parte superior del tobogán. (d) Exprese la energía total del sistema cuando la niña está en el punto de lanzamiento. (e) Exprese la energía total del sistema cuando la niña está en el aire $y_{\text{máx}}$ en términos de $g$ y $y_{\text{máx}}$. (f) A partir de los incisos (c) y (d), determine su rapidez inicial $v_i$ en términos de $g$ y $h$. (g) De los incisos (d), (e) y (f), determine su altura máxima en el aire $y_{\text{máx}}$ en términos de $h$ y del ángulo de lanzamiento $\theta$. (h) ¿Sus respuestas serían las mismas si el tobogán no tuviera fricción? Explique.
    \vspace*{9cm}
\end{itemize}

\subsection*{SECCIÓN 8.5 Potencia}

\begin{itemize}
    \item[\textbf{14.}] Un paracaidista de $80.0 \text{ kg}$ salta de un globo a una altura de $1000 \text{ m}$ y abre el paracaídas a una altitud de $200 \text{ m}$. (a) Si supone que la fuerza retardadora total sobre el paracaidista es constante en $3600 \text{ N}$ con el paracaídas abierto, encuentre la rapidez del paracaidista cuando aterriza en el suelo. (b) ¿Cree que el paracaidista se lesionará? Explique. (c) ¿A qué altura se debe abrir el paracaídas de modo que la rapidez final del paracaidista cuando golpee el suelo sea $5.00 \text{ m/s}$? (d) ¿Qué tan real es la suposición de que la fuerza retardadora total es constante? Explique.
    \vspace*{7cm}

    \item[\textbf{15.}] Usted ha pasado un largo día esquiando y está cansado. Está de pie en la parte superior de una colina, mirando hacia el albergue que está en la parte inferior de la colina. Está so cansado que quiere simplemente llegar a descansar y baja la ladera, sin empujar con sus bastones o hacer cualquier otra cosa para cambiar su movimiento. ¡Quiere que la gravedad haga todo el trabajo! Usted tiene la opción de dos senderos para llegar al albergue. Ambos senderos tienen el mismo coeficiente de fricción $\mu_k$. Además, ambos senderos representan la misma separación horizontal entre los puntos inicial y final. El sendero A tiene una bajada corta y empinada y luego un deslizamiento largo y plano hacia albergue. El sendero B tiene una pendiente descendente larga y ligera y luego un restante deslizamiento corto y plano al albergue. ¿Con cuál sendero llegará al albergue con la rapidez final más grande?
    \vspace*{6cm}

    \item[\textbf{16.}] El motor eléctrico de un tren a escala acelera al tren desde el reposo a $0.620 \text{ m/s}$ en $21.0 \text{ ms}$. La masa total del tren es $875 \text{ g}$. (a) Encuentre la potencia mínima entregada por la transmisión eléctrica de los rieles de metal durante la aceleración. (b) ¿Por qué la potencia es mínima?
    \vspace*{6cm}
\end{itemize}

\newpage

\begin{itemize}
    \item[\textbf{17.}] Una lámpara con eficiencia energética, que toma $28.0 \text{ W}$ de potencia, produce el mismo nivel de brillantez que una lámpara convencional que funciona a una potencia de $100 \text{ W}$. El tiempo de vida de la lámpara con eficiencia energética es $10\,000 \text{ h}$ y su precio de compra es $4.50$ dólares, mientras que la lámpara convencional tiene un tiempo de vida de $750 \text{ h}$ y cuesta $0.420$ dólares por lámpara. Determine el ahorro total que se obtiene al usar una lámpara con eficiencia energética durante su tiempo de vida, en oposición a usar lámparas convencionales durante el mismo intervalo de tiempo. Suponga un costo de energía de $0.200$ dólares por kilowatt hora.
    \vspace*{7cm}

    \item[\textbf{18.}] Un modelo viejo de automóvil acelera de $0$ a rapidez $v$ en un intervalo de tiempo $\Delta t$. Un automóvil deportivo nuevo y más potente acelera de $0$ a $2v$ en el mismo periodo. Suponiendo que la energía que proviene del motor aparece solo como energía cinética de los autos, compare la potencia de los dos automóviles.
    \vspace*{6cm}

    \item[\textbf{19.}] Haga la estimación de un orden de magnitud de la potencia que aporta el motor de un automóvil para acelerar el auto a rapidez de autopista. En su solución, establezca las cantidades que toma como datos y los valores que mide o estima para ellas. La masa del vehículo se proporciona en el manual del propietario.
    \vspace*{6cm}

    \item[\textbf{20.}] Habrá una caminata 5K que tendrá lugar en su ciudad. Mientras habla con su abuela, que usa una motoneta eléctrica para moverse, ella le dice que le gustaría acompañarle con su motoneta mientras usted camina la distancia de $5.00 \text{ km}$. El manual que venía en la motoneta indica que con batería completamente cargada puede proporcionar $120 \text{ Wh}$ de energía antes de que se agote. Para prepararse para la carrera, usted va para una ``unidad de prueba'': comenzando con una batería completamente cargada, su abuela va a su lado mientras camina $5.00 \text{ km}$ en terreno plano. Al final de la caminata, el indicador de uso de la batería muestra que permanece $40.0\%$ de la carga original en la batería. Usted también sabe que el peso combinado de la motoneta y su abuela es de $890 \text{ N}$. Unos días más tarde, lleno de confianza de que la batería tiene suficiente energía, usted y su abuela se dirigen hacia el evento 5K no es en terreno plano, sino que es cuesta arriba, terminando en un punto más alto que la línea de salida. Un funcionario de la carrera le dice que la cantidad total de desplazamiento vertical en la ruta es de $150 \text{ m}$. ¿Podría acompañarle su abuela en la caminata, o se quedará varada cuando la batería se agote? Suponga que la única diferencia entre su prueba y el evento real es el desplazamiento vertical.
    \vspace*{7cm}
\end{itemize}

\newpage

\begin{itemize}
    \item[\textbf{21.}] Para ahorrar energía, andar en bicicleta y caminar son medios de transporte mucho más eficientes que viajar en automóvil. Por ejemplo, al pedalear a $10.0 \text{ mi/h}$, un ciclista utiliza la energía del alimento a un ritmo de alrededor de $400 \text{ kcal/h}$ arriba de lo que usaría si simplemente estuviese sentado. (En fisiología del ejercicio, la potencia suele medirse en kcal/h en lugar de watts. Aquí $1 \text{ kcal} = 1 \text{ caloría de nutrición} = 4186 \text{ J}$.) Caminar a $3.00 \text{ mi/h}$ requiere unas $220 \text{ kcal/h}$. Además, comparar estos valores con el consumo de energía necesario para viajar en auto. La gasolina rinde aproximadamente $1.30 \times 10^4 \text{ kcal}$ por galón. Encuentre el combustible ahorrado en millas equivalentes por galón para (a) una persona caminando y (b) en bicicleta.
    \vspace*{7cm}

    \item[\textbf{22.}] Por convención la energía se mide en calorías, así como en joules. Una caloría en nutrición es una kilocaloría, que se define como $1 \text{ kcal} = 4\,186 \text{ J}$. Metabolizar $1 \text{ g}$ de grasa puede liberar $9.00 \text{ kcal}$. Un estudiante decide intentar perder peso mediante ejercicio. Planea subir y bajar corriendo las escaleras de un estadio de futbol tan rápido como pueda y tantas veces como sea necesario. Para evaluar el programa, suponga que sube un tramo de 80 escalones, cada uno de $0.150 \text{ m}$ de alto, en $65.0 \text{ s}$. Por simplicidad, desprecie la energía que usa al bajar (que es pequeña). Suponga que una eficiencia típica para músculos humanos es de $20.0\%$. Esta afirmación significa que, cuando su cuerpo convierte $100 \text{ J}$ metabolizando grasa, $20 \text{ J}$ realizan trabajo mecánico (en este caso, subir escaleras). El resto va a energía interna adicional. Suponga que la masa del estudiante es de $75.0 \text{ kg}$. (a) ¿Cuántas veces debe correr el tramo de escaleras para perder $1 \text{ kg}$ de grasa? (b) ¿Cuál es su potencia desarrollada promedio, en watts y en caballos de fuerza, mientras sube corriendo las escaleras? (c) ¿Esta actividad en sí misma es una forma práctica para perder peso?
    \vspace*{8cm}
\end{itemize}

\newpage

\section*{PROBLEMAS ADICIONALES}

\begin{itemize}
    \item[\textbf{23.}] Un bloque de masa $m = 200 \text{ g}$ se libera desde el reposo en el punto \textcircled{A} a lo largo del diámetro horizontal en el interior de un tazón hemisférico con radio $R = 30.0 \text{ cm}$ y la superficie del tazón es rugosa (figura P8.23). La rapidez del bloque en \textcircled{B} es $1.50 \text{ m/s}$. (a) ¿Cuál es su energía cinética en \textcircled{B}? (b) ¿Cuánta energía mecánica se transforma en energía interna?
    \vspace*{6cm}

    \item[\textbf{24.}] Haga una estimación de orden de magnitud de su potencia de salida cuando sube las escaleras. En su solución, indique las cantidades físicas que quiere tomar como datos y los valores que mide o estima para ellos. ¿Considera su potencia máxima o su energía sostenible?
    \vspace*{6cm}
\end{itemize}

\newpage

\begin{itemize}
    \item[\textbf{25.}] Usted está trabajando con un equipo que está diseñando una nueva montaña rusa del parque de atracciones para un parque temático importante. Usted está presente en la prueba de la montaña, en la que se envía un vagón de $250 \text{ kg}$ a lo largo de todo el trayecto. Cerca del final del juego, el vagón está casi en reposo en la parte superior de una pista de $110 \text{ m}$ de altura. Luego, entra en una sección final, rodando por una ondulante colina a nivel del suelo. La longitud total de la pista para esta sección final de la parte superior a la tierra es de $250 \text{ m}$. En los primeros $230 \text{ m}$, actúa una fuerza de fricción constante de $50.0 \text{ N}$ que actúa mediante unos frenos controlados por computadora. A nivel del suelo, la computadora aumenta la fuerza de fricción a un valor requerido para que la rapidez se reduzca a cero justo cuando el carro llega al punto de la pista que es la salida de pasajeros. (a) Determine la fuerza de fricción constante requerida de los últimos $20 \text{ m}$ para el vagón de prueba vacío. (b) Encuentre la rapidez máxima alcanzada por el vagón durante la sección final de pista de $250 \text{ m}$. (c) Su supervisor del equipo le pide determinar las respuestas a los incisos (a) y (b) para un vagón completamente cargado con un límite superior de masa de $450 \text{ kilogramos}$ por pasajero. Encuentre estos nuevos valores. (d) La fuerza de fricción requerida en el inciso (c) está bien dentro de los límites de diseño. Sin embargo, bajo la rapidez máxima, está muy por debajo de la de los principales juegos actuales, por lo que no puede hacer la torre más alta sobre el suelo, por lo que decide incluir un rasgo donde parte de la pista va \textit{bajo la tierra}. Determine la profundidad a la que debe estar la parte subterránea del trayecto para aumentar la rapidez máxima a $55.0 \text{ m/s}$. Suponga que la longitud total de la primera parte de la pista permanece a $230 \text{ m}$ y que la longitud de la pista de la parte superior a la subterránea en el punto más bajo es $150 \text{ m}$. La misma fuerza de fricción de $50.0 \text{ N}$ actúa en toda la sección de $230 \text{ m}$ de la pista. (e) ¿Es factible la construcción del inciso (d)?
    \vspace*{8cm}

    \item[\textbf{26.}] \textbf{Problema de repaso.} Como se muestra en la figura P8.26, una cuerda ligera que no se estira cambia de horizontal a vertical a medida que pasa sobre el borde de una mesa. La cuerda conecta un bloque de masa $m_1$ de $3.50 \text{ kg}$, al principio en reposo sobre la mesa horizontal de una altura $h = 1.20 \text{ m}$ arriba del suelo, un bloque que cuelga de masa $m_2$ de $1.90 \text{ kg}$ al principio a $0.900 \text{ m}$ sobre el suelo. Ni la superficie de la mesa ni su borde ejercen una fuerza de fricción cinética. Los bloques comienzan a moverse a partir del reposo. El bloque deslizante de masa $m_1$, se proyecta horizontalmente después de alcanzar el borde de la mesa. El bloque que cuelga de masa $m_2$ se para sin rebotar cuando golpea el piso. Considere a los dos bloques más la tierra como el sistema. (a) Encuentre la rapidez con la que $m_1$ deja el borde de la mesa. (b) Encuentre la rapidez de impacto de $m_1$ con el piso. (c) ¿Cuál es la longitud más corta de la cuerda que no está tensa mientras $m_1$ está en vuelo? (d) ¿Es la energía del sistema cuando se libera desde el reposo igual a la energía del sistema antes de que $m_1$ pegue en el suelo? (e) ¿Por qué o por qué no?
    \vspace*{8cm}

    \item[\textbf{27.}] Considere el sistema bloque-resorte-superficie en la parte (B) del ejemplo 8.6. (a) Usando la aproximación de energía encuentre la posición $x$ del bloque en el que rapidez es un máximo? (b) En la sección ¿Qué pasaría si?, de dicho ejemplo, se exploraron los efectos de una fuerza de fricción aumentada de $10.0 \text{ N}$. ¿En qué posición del bloque se presenta su rapidez máxima en esta situación?
    \vspace*{6cm}
\end{itemize}

\newpage

\begin{itemize}
    \item[\textbf{28.}] \textbf{¿Por qué no es posible la siguiente situación?} Una lanzadora de softbol tiene una extraña técnica: comienza con su mano en reposo en el punto más alto que puede alcanzar y luego gira rápidamente su brazo hacia atrás para que la bola se mueva a través de una trayectoria semicircular. Libera la bola cuando la mano llega a la parte inferior de la trayectoria. La lanzadora mantiene un componente de la fuerza sobre la bola de $0.180 \text{ kg}$ de magnitud constante $12.0 \text{ N}$ en la dirección del movimiento alrededor de la trayectoria completa. Cuando la pelota llega a la parte inferior de la trayectoria, deja su mano con una rapidez de $25.0 \text{ m/s}$.
    \vspace*{6cm}

    \item[\textbf{29.}] Jonathan anda en bicicleta y se encuentra con una colina de $7.30 \text{ m}$ de altura. En la base de la colina, viaja a $6.00 \text{ m/s}$. Cuando llega a la cima de la colina, está viajando a $1.00 \text{ m/s}$. Jonathan y su bicicleta juntos tienen una masa de $85.0 \text{ kg}$. Desprecie la fricción en el mecanismo de la bicicleta y entre los neumáticos de la bicicleta y la carretera. (a) ¿Cuál es el trabajo externo total hecho sobre el sistema de Jonathan y la bicicleta entre el tiempo que empieza a subir la colina y el tiempo que llega a la cima? (b) ¿Cuál es el cambio en energía potencial almacenada en el cuerpo de Jonathan durante este proceso? (c) ¿Cuánto trabajo hace Jonathan al pedalear la bicicleta dentro del sistema Jonathan-bicicleta-Tierra durante este proceso?
    \vspace*{6cm}

    \item[\textbf{30.}] Jonathan anda en bicicleta y se encuentra con una colina de altura $h$. En la base de la colina, está viajando con una rapidez $v_i$. Cuando llega a la cima de la colina, está viajando con una rapidez $v_f$. Jonathan y su bicicleta juntos tienen una masa $m$. Desprecie la fricción en el mecanismo de la bicicleta y entre los neumáticos de la bicicleta y la carretera. (a) ¿Cuál es el trabajo externo total hecho sobre el sistema de Jonathan y la bicicleta entre el tiempo que empieza a subir la colina y el tiempo que llega a la cima? (b) ¿Cuál es el cambio en energía potencial almacenada en el cuerpo de Jonathan durante este proceso? (c) ¿Cuánto trabajo hace Jonathan al pedalear la bicicleta dentro del sistema Jonathan-bicicleta-Tierra durante este proceso?
    \vspace*{6cm}

    \item[\textbf{31.}] Conforme el conductor pisa el pedal del acelerador, un automóvil de $1\,160 \text{ kg}$ de masa se acelera desde el reposo. Durante los primeros segundos de movimiento, el automóvil aumenta su aceleración con el tiempo de acuerdo con la expresión $a = 1.16t - 0.210t^2 + 0.240t^3$ donde $t$ está en segundos y $a$ está en $\text{m/s}^2$. (a) ¿Cuál es el cambio en la energía cinética del vehículo durante el intervalo de $t = 0$ a $t = 2.50 \text{ s}$? (b) ¿Cuál es la potencia promedio mínima del motor en este intervalo de tiempo? (c) ¿Por qué el inciso (b) se describe como el \textit{valor mínimo}?
    \vspace*{6cm}
\end{itemize}

\newpage

\begin{itemize}
    \item[\textbf{32.}] Mientras limpia un estacionamiento, un quitanieves empuja una pila cada vez más grande de nieve enfrente de él. Suponga que un automóvil que se mueve a través del aire se modela como un cilindro que empuja una pila creciente de nieve enfrente de él. El aire originalmente estacionario se pone en movimiento a la rapidez constante $v$ del cilindro, como se muestra en la figura P8.32. En un intervalo de tiempo $\Delta t$, un nuevo disco de aire de masa $\Delta m$ se debe mover una distancia $v\Delta t$ y por tanto se le debe dar una energía cinética $\frac{1}{2}(\Delta m)v^2$. Con el uso de este modelo, muestre que la pérdida de potencia del automóvil debida a la resistencia del aire es $\frac{1}{2}A \rho v^3$, donde $\rho$ es la densidad del aire. Compare este resultado con la expresión empírica $\frac{1}{2}DpAv^2$ para la fuerza resistiva.
    \vspace*{6cm}

    \item[\textbf{33.}] Haciendo caso omiso del peligro, un niño salta sobre una pila de colchones viejos para usarlos como un trampolín. Su movimiento entre dos puntos concretos se describe por la ecuación de conservación de energía: $\frac{1}{2}(46.0 \text{ kg})(2.40 \text{ m/s})^2 + (46.0 \text{ kg})(9.80 \text{ m/s}^2)(2.80 \text{ m} + x) = \frac{1}{2}(1.94 \times 10^4 \text{ N/m})x^2$. (a) Resuelva la ecuación para $x$. (b) Redacte el enunciado de un problema, incluyendo los datos para que esta ecuación sea la solución. (c) Agregue los dos valores de $x$ obtenidos en el inciso (a) y divida entre 2. (d) ¿Cuál es el significado del valor resultante en el inciso (c)?
    \vspace*{6cm}

    \item[\textbf{34.}] \textbf{Problema de repaso.} ¿Por qué no es posible la siguiente situación? Una nueva montaña rusa de alta rapidez se afirma que es tan segura que los pasajeros no necesitan llevar cinturones de seguridad o cualquier otro dispositivo de restricción. La montaña está diseñada con una sección circular vertical sobre el cual la montaña viaja en el interior del círculo para que los pasajeros estén boca abajo durante un corto intervalo de tiempo. El radio de la sección circular es $12.0 \text{ m}$ y la montaña rusa entra en la parte inferior de la sección circular con una rapidez de $22.0 \text{ m/s}$. El carro se mueve sin fricción en la pista y modele al carro como una partícula.
    \vspace*{6cm}

    \item[\textbf{35.}] Un resorte horizontal unido a una pared tiene una constante elástica $k = 850 \text{ N/m}$. Un bloque de masa $m = 1.00 \text{ kg}$ se une al resorte y descansa sobre una superficie horizontal sin fricción como en la figura P8.35. (a) El bloque se jala a una posición $x_i = 6.00 \text{ cm}$ desde la posición de equilibrio y se suelta. Encuentre la energía potencial elástica almacenada en el resorte cuando el bloque está $6.00 \text{ cm}$ de la posición de equilibrio y cuando el bloque pasa por la posición de equilibrio. (b) Encuentre la rapidez del bloque cuando pasa por el punto de equilibrio. (c) ¿Cuál es la rapidez del bloque cuando está en una posición $x/2 = 3.00 \text{ cm}$? (d) ¿Por qué no es la respuesta al inciso (c) la mitad de la respuesta del inciso (b)?
    \vspace*{6cm}
\end{itemize}

\newpage

\begin{itemize}
    \item[\textbf{36.}] Hace más de 2 300 años el maestro griego Aristóteles escribió el primer libro llamado \textit{Física}. Puesto en terminología más precisa, este pasaje es del final de su sección Eta: \textit{Sea P la potencia de un agente que causa movimiento; w, la carga movida; d, la distancia cubierta; y $\Delta t$, el intervalo de tiempo requerido.} Entonces (1) una potencia igual a $P$ en un intervalo de tiempo igual a $\Delta t$ moverá $w/2$ una distancia $d/2$ en el intervalo de tiempo $\Delta t/2$. Además, si (3) la potencia conocida $P$ mueve la carga dada $w$ una distancia $d/2$ en el intervalo de tiempo $\Delta t/2$, por tanto (4) $P/2$ moverá $w/2$ la distancia $d$ en el intervalo de tiempo dado $\Delta t$. (a) Demuestre que las proporciones de Aristóteles se incluyen en la ecuación $P\Delta t = bwd$, donde $b$ es una constante de proporcionalidad. (b) Demuestre que la teoría de movimiento del libro incluye esta parte de la teoría de Aristóteles como un caso especial. En particular, describa una situación en la que sea verdadera, deduzca la ecuación que represente las proporciones de Aristóteles y determine la constante de proporcionalidad.
    \vspace*{7cm}

    \item[\textbf{37.}] \textbf{Problema de repaso.} Como una broma, alguien equilibró una calabaza en el punto más alto de un gran silo. El silo tiene como remate un casco semiesférico que no tiene fricción cuando está mojado. La línea desde el centro de curvatura del casco hacia la calabaza forma un ángulo $\theta_i = 0^\circ$ con la vertical. En una noche lluviosa, mientras está de pie en las cercanías, un soplo de viento hace que la calabaza se comience a deslizar hacia abajo desde el reposo. La calabaza pierde contacto con el casco cuando la línea desde el centro del hemisferio hacia la calabaza forma cierto ángulo con la vertical. ¿Cuál es este ángulo?
    \vspace*{6cm}

    \item[\textbf{38.}] \textbf{Problema de repaso.} ¿Por qué no es posible la siguiente situación? Una atleta pone a prueba la fortaleza de sus manos al tener un asistente que cuelga pesos de su cinturón cuando se cuelga con sus manos de una barra horizontal. Cuando los pesos que cuelgan de su cinturón han aumentado hasta $80\%$ de su peso corporal, sus manos ya no pueden apoyarla y cae al piso. Frustrada por no cumplir su meta de fuerza en las manos, decide oscilar en un trapecio. El trapecio consiste en una barra suspendida por dos cuerdas paralelas, cada una de longitud $\ell$, permitiendo a los artistas oscilar en un arco circular vertical, comenzando a partir del reposo con las cuerdas en un ángulo $\theta_i = 60.0^\circ$ con respecto a la vertical. Cuando se balancea hacia adelante y hacia atrás varias veces en un arco circular, olvida su frustración relacionada con la prueba de fuerza en las manos. Suponga que el tamaño del cuerpo de la artista es pequeño en comparación con la longitud $\ell$ y que la resistencia del aire es despreciable.
    \vspace*{6cm}

    \item[\textbf{39.}] Un avión de $1.50 \times 10^4 \text{ kg}$ de masa hace un vuelo a nivel, moviéndose inicialmente a $60.0 \text{ m/s}$. La fuerza resistiva ejercida por el aire en el avión tiene una magnitud de $4.0 \times 10^5 \text{ N}$. Por la tercera ley de Newton, si los motores ejercen una fuerza en los gases de escape para expulsarlos de la parte posterior del motor, los gases de escape ejercen una fuerza en los motores en la dirección de desplazamiento del avión. Esta fuerza se denomina empuje, y el valor del empuje en esta situación es $7.50 \times 10^5 \text{ N}$. (a) ¿Es el trabajo realizado por los gases de escape en el avión durante un intervalo de tiempo igual al cambio en la energía cinética del avión? Explique. (b) Encuentre la rapidez del avión después de que ha viajado $5.0 \times 10^4 \text{ m}$.
    \vspace*{6cm}
\end{itemize}

\newpage

\section*{PROBLEMAS DE DESAFÍO}

\begin{itemize}
    \item[\textbf{40.}] Un péndulo, que consta de una cuerda ligera de longitud $L$ y una esfera pequeña, se balancea en el plano vertical. La cuerda golpea una clavija ubicada a una distancia $d$ bajo el punto de suspensión (figura P8.40). (a) Demuestre que, si la esfera se libera desde una altura por abajo de la clavija, regresará a esta altura después de que la cuerda golpee la clavija. (b) Demuestre que, si el péndulo se libera desde la posición horizontal ($\theta = 90^\circ$) y se balancea en un círculo completo con centro en la clavija, el valor mínimo de $d$ debe ser $3L/5$.
    \vspace*{7cm}

    \item[\textbf{41.}] Una bola gira alrededor de un círculo vertical en el extremo de una cuerda. El otro extremo de la cuerda está fijo en el centro del círculo. Suponiendo que la energía total del sistema bola-Tierra permanece constante, demuestre que la tensión en la cuerda en la parte baja es mayor que la tensión en la parte alta por seis veces el peso de la bola.
    \vspace*{6cm}

    \item[\textbf{42.}] Usted está trabajando en el centro de distribución de un gran sitio de compras en línea. Se están haciendo esfuerzos para aumentar el número de paquetes por unidad de tiempo que se están cargando en una banda transportadora para montarlos a camiones en espera, pero el motor impulsor de la banda transportadora tiene dificultad para mantener las crecientes demandas. Su supervisor le ha pedido determinar los requisitos para que un motor nuevo pueda proporcionar la energía suficiente para mantener la cinta transportadora en movimiento continuo sometido a la creciente rapidez de carga. Se le da la información siguiente: la meta del diseño es tener $50.0 \text{ kg}$ de paquetes cargados en la banda en diferentes posiciones en una cantidad promedio de $5.00$ paquetes por segundo. La banda se mueve con una rapidez horizontal de $1.35 \text{ m/s}$. Personas en diferentes lugares a lo largo de la banda colocan el paquete en la banda de modo que está inicialmente en reposo con respecto al suelo de la construcción en la banda. Su tarea es determinar la potencia mínima que debe tener el motor impulsor para acelerar estos paquetes y mantener la banda moviéndose con rapidez constante.
    \vspace*{7cm}

    \item[\textbf{43.}] Considere la colisión bloque-resorte analizada en el ejemplo 8.8. (a) En el inciso (B), para la situación en que la superficie ejerce una fuerza de fricción en el bloque, demuestre que el bloque nunca llega de regreso a $x = 0$. (b) ¿Cuál es el valor máximo del coeficiente de fricción que permitiría al bloque regresar a $x = 0$?
    \vspace*{6cm}
\end{itemize}

\newpage

\begin{itemize}
    \item[\textbf{44.}] Desde el reposo, una persona de $64.0 \text{ kg}$ salta en \textit{bungee} desde un globo de aire caliente a $65.0 \text{ m}$ sobre el suelo. La cuerda \textit{bungee} tiene masa despreciable y longitud sin estirar de $25.8 \text{ m}$. Un extremo se amarra a la canasta del globo aerostático y el otro extremo a un arnés alrededor del cuerpo de la persona. La cuerda se modela como un resorte que obedece a la ley de Hooke con una constante elástica de $81.0 \text{ N/m}$, y el cuerpo de la persona se modela como partícula. El globo no se mueve. (a) Exprese la energía potencial gravitacional del sistema persona-Tierra como función de la altura variable $y$ de la persona sobre el suelo. (b) Exprese la energía potencial elástica de la cuerda como función de $y$. (c) Exprese la energía potencial total del sistema persona-Tierra como función de $y$. (d) Trace una gráfica de energías gravitacional, elástica y potencial total como funciones de $y$. (e) Suponga que la resistencia del aire es despreciable. Determine la altura mínima de la persona sobre el suelo durante su caída. (f) ¿La gráfica de energía potencial muestra alguna posición de equilibrio? Si es así, ¿a qué elevaciones? ¿Son estables o inestables? (g) Determine la rapidez máxima del saltador.
    \vspace*{9cm}

    \item[\textbf{45.}] \textbf{Problema de repaso.} Un tablero uniforme de longitud $L$ se desliza a lo largo de un plano horizontal uniforme (sin fricción), como se muestra en la figura P8.45a. Después el tablero se desliza a través de la frontera con una superficie horizontal rugosa. El coeficiente de fricción cinética entre el tablero y la segunda superficie es $\mu_k$. (a) Encuentre la aceleración del tablero cuando su extremo frontal recorre una distancia $x$ más allá de la frontera. (b) El tablero se detiene en el momento en que su extremo posterior llega a la frontera, como se muestra en la figura P8.45b. Encuentre la rapidez inicial $v$ del tablero.
    \vspace*{6cm}

    \item[\textbf{46.}] Una cadena uniforme de $8.00 \text{ m}$ de longitud inicialmente yace estirada sobre una mesa horizontal. (a) Suponiendo que el coeficiente de fricción estática entre la cadena y la mesa es $0.600$, demuestre que la cadena comenzará a deslizarse de la mesa si al menos $3.00 \text{ m}$ de ella cuelgan sobre el borde de la mesa. (b) Determine la rapidez de la cadena cuando el extremo de la cadena cuando el último eslabón deja la mesa, teniendo en cuenta que el coeficiente de fricción cinética entre la cadena y la mesa es $0.400$.
    \vspace*{7cm}

    \item[\textbf{47.}] \textbf{¿Qué pasaría si?} Considere la montaña rusa que se describe en el problema 34. Debido a que hay algo de fricción entre el carro y la pista, el carro entra en la sección circular con una rapidez de $15.0 \text{ m/s}$ en lugar de los $22.0 \text{ m/s}$. ¿Esta situación es más o menos peligrosa para los pasajeros que la del problema 34? Suponga que la sección circular aún está sin fricción.
    \vspace*{6cm}
\end{itemize}

\end{document}
